\section{Impl\'ementation : chronologie et pi\`eges}
\label{sec:implementation}

\begin{resume}
Le chantier a suivi dix t\^aches, chacune \'ecrite en cycle test d'abord :
un test rouge cibl\'e, une impl\'ementation minimale, une v\'erification, un
commit. Les fondations (chronom\'etrage, \'evaluateur injectable, cache par
snapshots, r\'eservations RAII, pool persistant) pr\'ec\`edent la collecte
parall\`ele, le coordinateur, le stress, les bancs et enfin les mesures. Les
deux incidents marquants sont un contrat de dimensions silencieux d\'ecouvert
par le banc self-play, et un test d'\'equivalence du pool chaud dont la
premi\`ere version comparait la mauvaise r\'ef\'erence. Ce chapitre raconte
l'ordre r\'eel des d\'ecouvertes, avec les commits correspondants.
\end{resume}

\subsection{T\^ache 1 : la r\'ef\'erence et le chronom\'etrage}
\label{sec:t1}

Premi\`ere t\^ache, premi\`ere discipline : mesurer avant de modifier. La
r\'ef\'erence a \'et\'e captur\'ee avec le binaire de d\'epart, puis dix
phases de recherche ont \'et\'e instrument\'ees : s\'election, pr\'eparation du
tenseur et de la cl\'e, lecture et \'ecriture de table, attente de verrou,
copie du plateau, assemblage, inf\'erence, expansion, remont\'ees, attente
des workers. Le chronom\'etrage est d\'esactiv\'e par d\'efaut et sans co\^ut
mesurable dans ce mode : quand le pointeur de mesure est nul, aucune horloge
n'est lue.

La campagne de r\'ef\'erence a donn\'e la r\'epartition qui a orient\'e tout
le reste : l'inf\'erence repr\'esente environ 98\,\% du temps. C'est cette
mesure, plus que toute intuition, qui a fix\'e l'objectif r\'ealiste du
chantier multic\oe{}ur : parall\'eliser au mieux 2\,\% du temps, et r\'eduire
les lots partiels. Commit \code{4b4f92f}.

\subsection{T\^ache 2 : l'\'evaluateur injectable et le banc self-play}
\label{sec:t2}

Le MCTS ne connaissait que \code{ONNXEvaluator}. Une interface
\code{Evaluator} a \'et\'e extraite, avec une seule m\'ethode
\code{evaluate\_batch} ; \code{ONNXEvaluator} en h\'erite sans changer de
comportement, et un \code{ControlledEvaluator} de test permet de simuler des
\'echecs, des lots de tailles choisies et des blocages. En parall\`ele, les
copies de \code{Chessboard} ont \'et\'e verrouill\'ees par des tests sur des
histoires l\'egales concr\`etes : roque, prise en passant, promotion,
r\'ep\'etition, avec comparaison de la FEN, du Zobrist, du tenseur et des
historiques apr\`es copie, coup jou\'e et annul\'e.

Ce travail a aussi produit le banc \code{selfplay\_phase\_bench}, qui rejoue la
structure du self-play partag\'e : huit racines priv\'ees, cent tours de
collecte sous OpenMP, puis une \'evaluation et des remont\'ees
s\'equentielles. C'est ce banc qui a servi plus tard de d\'etecteur de
r\'egression. Deux accrocs mineurs y ont \'et\'e corrig\'es : une condition de
validation invers\'ee dans le nouvel outil en ligne de commande, et la
d\'ecouverte que l'ex\'ecutable autonome a besoin du r\'epertoire
\code{torch/lib} dans son \code{PATH} pour trouver \code{cublasLt}. Commit
\code{fe63dbe}.

\subsection{T\^ache 3 : le cache par snapshots}
\label{sec:t3}

La table de transposition est pass\'ee derri\`ere une classe
\code{EvaluationCache}. Le point de bascule est le type de retour : le
\code{probe} ne rend plus un pointeur, mais une copie. Pour \'eliminer le
verrou global, l'index est d\'ecoup\'e en au plus 4096 bandes, d\'eriv\'ees
de l'index physique. Les tests de collision forcent deux writers et deux
readers sur des cl\'es qui ciblent la m\^eme entr\'ee, et v\'erifient qu'un
snapshot ne m\'elange jamais deux politiques. La migration a \'et\'e men\'ee
en deux temps : remplacer les types d'abord, compiler rouge, puis corriger
chaque usage list\'e par le compilateur.

Le surco\^ut mesur\'e face au r\'eseau est minuscule : environ
$0{,}27$ \`a $0{,}33$\,ms d'attente de verrou et $0{,}45$ \`a
$0{,}50$\,ms d'acc\`es \`a la table pour une recherche de 0,67 \`a
1,46 seconde. Sur le faux r\'eseau, o\`u ce co\^ut n'est plus masqu\'e, la
collecte passe de 0,79 \`a 1,14\,ms pour 100 vagues, soit environ 3,5
microsecondes ajout\'ees par vague. Commit \code{a0a182d}.

\subsection{T\^ache 4 : \'etats atomiques et r\'eservations RAII}
\label{sec:t4}

L'ancien bool\'een \code{is\_terminal} a disparu au profit de
l'\'enum\'eration d'\'etats d\'ecrite au chapitre~3, et le compteur
\code{n\_in\_flight} est devenu un entier atomique. La classe
\code{PathReservation} poss\`ede le chemin et le \code{Pending} \'eventuel.
Un test de mutation a \'et\'e pratiqu\'e : en neutralisant volontairement la
d\'ecr\'ementation du compteur, le test de fuite \'echoue comme pr\'evu, puis
le code correct est restaur\'e. La garde a tenu cinquante ex\'ecutions
concurrentes. Commit \code{a303d96}.

\subsection{T\^ache 5 : le pool de workers persistant}
\label{sec:t5}

\code{SearchExecutor} cr\'ee les threads au premier appel et les rendort
entre les vagues sur une variable de condition et un num\'ero de
g\'en\'eration. C'est le rendez-vous C++17 demand\'e par les contraintes : pas
de \code{std::jthread}, pas de \code{std::barrier}. Les tests v\'erifient que
les identit\'es de threads survivent aux appels (8, puis 2, puis 8), qu'une
exception de worker est relay\'ee apr\`es le retour de tous, et qu'un appel
ult\'erieur reste possible apr\`es une exception. Le pool ne grandit qu'au
repos. Commit \code{540e5ee}.

\subsection{T\^ache 6 : le collecteur parall\`ele}
\label{sec:t6}

Le fichier \code{mcts\_wave}\allowbreak\code{.cpp} contient la descente
concurrente. Chaque
worker a son \code{WorkerContext} (plateau complet, r\'esultats locaux,
chronom\^etre). La boucle de collecte r\'eserve la racine, s\'electionne avec
le PUCT existant, r\'eserve l'enfant imm\'ediatement, applique le coup sur le
plateau priv\'e et recommence. Les points d\'ecision sont les suivants.

\begin{lstlisting}[caption={Extrait de la boucle de collecte. Le terminal est
test\'e avant toute lecture de table ; le CAS vers \code{Pending} est la seule
garantie de propri\'et\'e ; un succ\`es de table publie les enfants et
poursuit la descente, comme le faisait \code{select\_leaf}.},
label={lst:collecte}]
// ... descente : etat du noeud, terminal par les regles du plateau ...
if (is_rule_terminal(context.board)) {
    if (state == NodeState::Pending) return collision();
    if (state == NodeState::Unexpanded) {
        if (!work.reservation.try_claim(node)) return collision();
        work.reservation.publish(NodeState::Terminal);
    }
    work.kind = LeafKind::Terminal;
    return work;
}
if (state == NodeState::Unexpanded) {
    if (!work.reservation.try_claim(node)) return collision();
    legal_moves = context.board.getLegalMoveIndices();
    if (legal_moves.empty()) {
        work.reservation.publish(NodeState::Terminal);
        return work;   // mat ou pat, valeur terminale
    }
    key = make_cache_key(context.board);
    probe = probe_tt(key, &context.timing);
    if (probe.status == TTProbeStatus::HIT) {
        children = make_children_from_probe(node, probe);
        node->children.swap(children);          // publication complete
        work.reservation.publish(NodeState::Expanded);
        continue;                                // meme descente
    }
    // vraie feuille : capturer tenseur, coups, cle, transferer la garde
    context.board.getAlphaZeroTensor(work.tensor);
    work.kind = LeafKind::Network;
    return work;
}
// ... noeud Expanded : selection PUCT, reserve de l'enfant, coup joue ...
\end{lstlisting}

Les permis bornent la vague : \code{slots = min(lot, restant)}, un CAS
d\'ecr\'emente \code{available}, une collision le rend, un r\'esultat final le
consomme. Le nombre de tentatives est plafonn\'e \`a $4 \times \text{slots} +
\text{workers}$. Les vectors de r\'esultats restent locaux \`a chaque worker
et sont fusionn\'es apr\`es le rendez-vous dans un ordre stable
\code{(worker, ordinal)}. Deux tests d\'ecisifs ont \'et\'e \'ecrits \`a ce
moment : la comparaison du tenseur captur\'e avec le tenseur obtenu en
rejouant r\'eellement l'historique, et la publication atomique depuis la
table avec un point de crochet avant publication. Cent ex\'ecutions
concurrentes ont valid\'e l'ensemble. Commit \code{2772087}.

\subsection{T\^ache 7 : le coordinateur et l'API}
\label{sec:t7}

\code{run\_search\_waves} orchestre les vagues : copie du plateau racine dans
les seuls contextes actifs, collecte, fusion, une inf\'erence si au moins une
feuille r\'eseau, validation des dimensions et des valeurs finies, puis
expansion et remont\'ees dans l'ordre. Les compteurs nouveaux
(\code{waves}, \code{leaf\_collisions}, \code{completed\_simulations}) sont
expos\'es c\^ot\'e Python. Le param\`etre \code{worker\_count} est ajout\'e
aux deux points d'entr\'ee publics, avec une valeur par d\'efaut de 1 qui
pr\'eserve tous les appelants existants ; une configuration
multi\oe{}ur avec un lot nul est refus\'ee avant toute mutation.

Le verrou de session est d\'esormais tenu pendant tout l'appel, pas par
simulation : les workers d\'etiennent des pointeurs bruts vers l'arbre, donc
\code{update\_root} ne peut pas d\'etruire la racine sous eux. Les tests
int\'egr\'es couvrent les workers 2, 4 et 8, les lots 1 et 8, les budgets 0,
1, 3, 8, 17 et 100, les racines terminales, les erreurs de dimensions de
l'\'evaluateur, l'immobilit\'e des workers pendant une inf\'erence bloqu\'ee,
et la s\'erialisation de \code{update\_root} par un \'evaluateur retenu.
Commit \code{3b42dd2}.

\subsection{T\^ache 8 : le stress et le self-play partag\'e}
\label{sec:t8}

Cette t\^ache ne changeait pas l'algorithme : elle cherchait \`a casser le
nouveau noyau. Trois positions, workers $\{2, 4, 8\}$, lots $\{1, 8\}$, cent
recherches par configuration, en alternant racines neuves et r\'eutilis\'ees,
puis en alternant $8 \to 2 \to 1 \to 4$ : 1\,800 recherches et 115\,200
simulations, invariants v\'erifi\'es apr\`es chaque appel. Le self-play
partag\'e a \'et\'e exerc\'e avec huit racines priv\'ees sur une instance
MCTS et une table minuscule, puis avec deux parties compl\`etes du manager.

\begin{piege}[Le contrat de dimensions, silencieux]
Le test \code{selfplay\_shared\_core} a termin\'e ses deux parties puis
\'echou\'e : le faux \'evaluateur recevait un vecteur d'entr\'ee plus grand que
\code{batch\_size} $\times 119 \times 64$. La cause \'etait un buffer
d'entr\'ee r\'eutilis\'e, redimensionn\'e une fois pour le nombre maximal de
parties, mais jamais ramen\'e \`a la taille r\'eelle du lot en cours de
traitement. Un vrai \'evaluateur ONNX ne se plaint pas de cette diff\'erence :
le contrat de dimensions n'\'etait v\'erifi\'e que par le faux r\'eseau, et
c'est pr\'ecis\'ement pour cela que le bug a \'et\'e vu.
\end{piege}

Deux corrections ont suivi : redimensionner le buffer juste avant chaque
appel dans les deux chemins du manager, et faire d\'ependre
\code{SelfPlayManager} de l'interface \code{Evaluator} plut\`o\`ut que de la
classe ONNX concr\`ete, ce qui \'etait la seule adaptation de production
n\'ecessaire. Deux tests de mutation ont \'et\'e pratiqu\'es : neutraliser la
remise \`a z\'ero de \code{Pending} doit faire \'echouer le test d'exception,
et neutraliser le CAS de propri\'et\'e doit faire \'echouer le test de
collision. Commits \code{ef327e2} et \code{9ad2dcb}.

\subsection{T\^ache 9 : les bancs comparables}
\label{sec:t9}

Les bancs existants mesuraient un arbre neuf par mesure, sans notion de
workers ni de pool. Ils ont gagn\'e quatre capacit\'es : un balayage de
\code{--worker-counts}, un mode \code{--warmup-pool} qui r\'eutilise l'objet
MCTS entre les mesures en remettant arbre et table \`a froid hors
chronom\'etrage, une jambe GPU pour le banc de puzzles avec
\code{--search-workers} et \code{--search-seconds}, et un sidecar de
contexte par CSV (hash du mod\`ele et du banc, budget, provider, workers,
crit\`ere fig\'e). La comparaison de campagnes a re\c{c}u son propre module,
\code{multicore\_comparison.py} : refus des doublons, des lignes manquantes et
des m\'etadonn\'ees incompatibles, bootstrap appari\'e \`a 20\,000 tirages,
paires discordantes, McNemar, et verdict formel. Une passe de durcissement a
ensuite ajout\'e le statut d'\'echec de campagne quand un puzzle est en
erreur, les bornes exactes des verdicts, et l'alternance de l'ordre des
workers entre passages. Commits \code{7bde9e3}, \code{7cef664} et
\code{9840fa8}.

\begin{piege}[Le premier test d'\'equivalence du pool chaud \'etait faux]
Pour v\'erifier qu'un \code{clear\_evaluation\_cache()} ram\`ene bien un
objet au m\^eme \'etat qu'un objet neuf, un test comparait la deuxi\`eme
recherche de l'objet r\'eutilis\'e \`a la deuxi\`eme recherche d'un objet
neuf. Il a \'echou\'e proprement, et pour la bonne raison : la deuxi\`eme
recherche de l'objet neuf b\'en\'eficiait de sa propre table d\'ej\`a
remplie, avec dix succ\`es de table contre z\'ero apr\`es le clear. La
r\'ef\'erence correcte est la \emph{premi\`ere} recherche d'un objet neuf.
Une fois corrig\'e, politiques et compteurs sont identiques au bit pr\`es,
ce qui valide \`a la fois le clear et la conception du banc.
\end{piege}

\subsection{T\^ache 10 : mesures et activation}
\label{sec:t10}

La derni\`ere t\^ache n'ajoute pas de m\'ecanisme : elle ex\'ecute les
campagnes, compare l'ancien et le nouveau binaire, mesure la qualit\'e,
r\'edige le rapport de r\'esultats et active le r\'eglage UCI. Le
chapitre~6 d\'etaille toutes les mesures. Commit \code{412b5d4} pour
l'activation, \code{0aaeeb1} pour les cases du plan.
